Kurs:Algebraische Kurven (Osnabrück 2025-2026)/Arbeitsblatt 21/latex
Kurs:Algebraische Kurven (Osnabrück 2025-2026)/Arbeitsblatt 21/latex

\setcounter{section}{21}


  \zwischenueberschrift{Übungsaufgaben}


\inputaufgabe

{}

{


Es sei


\mavergleichskette

{\vergleichskette

{ M
}

{ \subseteq }{\N
}

{  }{
}

{    }{
}

{   }{
}

}
{}{}{}
ein
\definitionsverweis {numerisches Monoid}{}{,}
das von teilerfremden Erzeugern erzeugt werde, es sei

\mathl{K[M]}{} der
\definitionsverweis {Monoidring}{}{}
zu $M$ über einem Körper $K$ und es sei


\mavergleichskettedisp

{\vergleichskette

{ R
}

{ =} { K[M]_{\mathfrak m}
}

{ } {
}

{ } {
}

{ } {
}

}
{}{}{}
die
\definitionsverweis {Lokalisierung}{}{}
am
\definitionsverweis {maximalen Ideal}{}{}


\mavergleichskette

{\vergleichskette

{ {\mathfrak m}
}

{ = }{ K[ M_+]
}

{ = }{ \langle  T^m ,\, m \in M  \rangle
}

{    }{
}

{   }{
}

}
{}{}{.}
Zeige, dass $R$ allein im Fall


\mavergleichskette

{\vergleichskette

{ M
}

{ = }{ \N
}

{  }{
}

{    }{
}

{   }{
}

}
{}{}{}
ein
\definitionsverweis {diskreter Bewertungsring}{}{}
ist.


}

{} {}


\inputaufgabe

{}

{


Es sei $R$ ein
\definitionsverweis {diskreter Bewertungsring}{}{}
mit
\definitionsverweis {maximalem Ideal}{}{}


\mavergleichskette

{\vergleichskette

{ {\mathfrak m}
}

{ = }{ (p)
}

{  }{
}

{    }{
}

{   }{
}

}
{}{}{.}
Zeige, dass die
\definitionsverweis {Ordnung}{}{}
\maabbeledisp {} {R \setminus \{0\} } { \N
} { f } { \operatorname{ord} \,  \left(   f   \right)
} {,}
folgende Eigenschaften besitzt.
\aufzaehlungvier{

\mavergleichskette

{\vergleichskette

{ \operatorname{ord} \,  \left(  fg  \right)
}

{ = }{ \operatorname{ord} \,  \left(   f   \right) + \operatorname{ord} \,  \left(   g   \right)
}

{  }{
}

{    }{
}

{   }{
}

}
{}{}{.}
}{

\mavergleichskette

{\vergleichskette

{ \operatorname{ord} \,  \left(  f+g  \right)
}

{ \geq }{ \operatorname{min} \left(  \operatorname{ord} \,  \left(   f   \right)  ,\,   \operatorname{ord} \,  \left(   g   \right)  \right)
}

{  }{
}

{    }{
}

{   }{
}

}
{}{}{.}
}{Es ist


\mavergleichskette

{\vergleichskette

{f
}

{ \in }{ {\mathfrak m}
}

{  }{
}

{    }{
}

{   }{
}

}
{}{}{}
genau dann, wenn


\mavergleichskette

{\vergleichskette

{ \operatorname{ord} \,  \left(   f   \right)
}

{ \geq }{ 1
}

{  }{
}

{    }{
}

{   }{
}

}
{}{}{}
ist.
}{Es ist


\mavergleichskette

{\vergleichskette

{f
}

{ \in }{ R ^{\times}
}

{  }{
}

{    }{
}

{   }{
}

}
{}{}{}
genau dann, wenn


\mavergleichskette

{\vergleichskette

{ \operatorname{ord} \,  \left(   f   \right)
}

{ = }{ 0
}

{  }{
}

{    }{
}

{   }{
}

}
{}{}{}
ist.
}


}

{} {}


\inputaufgabegibtloesung

{}

{


Es sei $R$ ein
\definitionsverweis {diskreter Bewertungsring}{}{}
mit
\definitionsverweis {Quotientenkörper}{}{}
$Q$. Zeige, dass es keinen echten Zwischenring zwischen $R$ und $Q$ gibt.


}

{} {}


\inputaufgabe

{}

{


Es sei $R$ ein
\definitionsverweis {diskreter Bewertungsring}{}{}
mit
\definitionsverweis {Quotientenkörper}{}{}
$Q$. Charakterisiere die
\definitionsverweis {endlich erzeugten}{}{}
$R$-\definitionsverweis {Untermoduln}{}{}
von $Q$. Auf welche Form kann man ein Erzeugendensystem bringen?


}

{} {}


\inputaufgabe

{}

{


Es sei $R$ ein
\definitionsverweis {diskreter Bewertungsring}{}{.}
Definiere zu einem Element

\mathbed {q \in Q(R)} {}

 {q \neq 0} {}

 {} {} {} {,} die Ordnung


\mavergleichskettedisp

{\vergleichskette

{ \operatorname{ord} \,  \left(   q   \right)
}

{ \in} { \Z
}

{ } {
}

{ } {
}

{ } {
}

}
{}{}{.}
Dabei soll die Definition mit der
\definitionsverweis {Ordnung}{}{}
für Elemente aus $R$ übereinstimmen und einen Gruppenhomomorphismus
\maabb {} {Q(R) \setminus \{0\} } { \Z
} {}
definieren. Was ist der Kern dieses Homomorphismus?


}

{} {}


\inputaufgabe

{}

{


Es sei


\mavergleichskette

{\vergleichskette

{ V
}

{ = }{ V \left(  x^2+y^2-1  \right)
}

{ \subseteq }{ {\mathbb A}^{2}_{ K }
}

{    }{
}

{   }{
}

}
{}{}{}
der \stichwort {Einheitskreis} {} über einem
\definitionsverweis {Körper}{}{}
$K$ und es sei


\mavergleichskette

{\vergleichskette

{ P
}

{ = }{ (a,b)
}

{ \in }{ V
}

{    }{
}

{   }{
}

}
{}{}{}
ein Punkt.
\aufzaehlungvierabc{Zeige, dass der
\definitionsverweis {lokale Ring}{}{}
$R$ von $V$ im Punkt $P$ ein
\definitionsverweis {diskreter Bewertungsring}{}{}
ist.
}{ Folgere, dass der
\definitionsverweis {Koordinatenring}{}{}


\mathl{K[X,Y]/ \left(  X^2+Y^2-1  \right)}{}
\definitionsverweis {normal}{}{}
ist
\zusatzklammer {man kann $K$
\definitionsverweis {algebraisch abgeschlossen}{}{}
annehmen} {} {.}
}{Zeige, dass

\mathl{K[X,Y]/ \left(  X^2+Y^2-1  \right)}{} nicht
\definitionsverweis {faktoriell}{}{}
ist.
}{Bestimme die
\definitionsverweis {Ordnung}{}{}
von $X$ und von $Y-1$ im lokalen Ring zum Punkt

\mathl{(0,1)}{.}
}


}

{} {}


\inputaufgabe

{}

{


Es sei $R$ ein
\definitionsverweis {diskreter Bewertungsring}{}{}
und sei


\mavergleichskette

{\vergleichskette

{ {\mathfrak m}
}

{ = }{ ( \pi)
}

{  }{
}

{    }{
}

{   }{
}

}
{}{}{.}
Es sei


\mavergleichskette

{\vergleichskette

{K
}

{ = }{ R/ (\pi)
}

{  }{
}

{    }{
}

{   }{
}

}
{}{}{}
der
\definitionsverweis {Restklassenkörper}{}{}
von $R$. Zeige, dass es für jedes


\mavergleichskette

{\vergleichskette

{ n
}

{ \in }{ \N
}

{  }{
}

{    }{
}

{   }{
}

}
{}{}{}
einen
$R$-\definitionsverweis {Modulisomorphismus}{}{}
\maabbdisp {} { (\pi^ n)/( \pi^ {n+1} )} { K
} {}
gibt.


}

{} {}


\inputaufgabegibtloesung

{}

{


Es sei $K$ ein
\definitionsverweis {Körper}{}{}
und sei
\maabbdisp {\nu} {(K^\times, \cdot,1)} { (\Z,+,0)
} {}
ein
\definitionsverweis {surjektiver}{}{}
\definitionsverweis {Gruppenhomomorphismus}{}{}
mit


\mavergleichskette

{\vergleichskette

{ \nu(f+g)
}

{ \geq }{ \min\{ \nu(f) , \nu(g)\}
}

{  }{
}

{    }{
}

{   }{
}

}
{}{}{}
für alle


\mavergleichskette

{\vergleichskette

{ f,g
}

{ \in }{ K ^{\times}
}

{  }{
}

{    }{
}

{   }{
}

}
{}{}{.}
Zeige, dass


\mavergleichskettedisp

{\vergleichskette

{R
}

{ =} { \left\{ f \in K^\times   \mid  \nu(f) \geq 0         \right\}   \cup \{0\}
}

{ } {
}

{ } {
}

{ } {
}

}
{}{}{}
ein
\definitionsverweis {diskreter Bewertungsring}{}{}
ist.


}

{} {}


\inputaufgabe

{}

{


Es sei


\mathbed {f \in {\mathbb C}[X]} {}

 {f \neq 0} {}

 {} {} {} {,}
und


\mavergleichskette

{\vergleichskette

{ a
}

{ \in }{ {\mathbb C}
}

{  }{
}

{    }{
}

{   }{
}

}
{}{}{.}
Zeige, dass die folgenden \anfuehrung{Ordnungen}{} von $f$ an der Stelle $a$ übereinstimmen.
\aufzaehlungdrei{Die Verschwindungsordnung von $f$ an der Stelle $a$, also die minimale Ordnung einer Ableitung mit


\mavergleichskette

{\vergleichskette

{ f^{(k)}(a)
}

{ \neq }{ 0
}

{  }{
}

{    }{
}

{   }{
}

}
{}{}{.}
}{Der Exponent des Linearfaktors

\mathl{X-a}{} in der Zerlegung von $f$ in irreduzible Polynome.
}{Die
\definitionsverweis {Ordnung}{}{}
von $f$ an der Lokalisierung

\mathl{{\mathbb C}[X]_{(X-a)}}{}  von

\mathl{{\mathbb C}[X]}{} am maximalen Ideal

\mathl{(X-a)}{.}
}


}

{} {}


Die vorstehende Aufgabe kann man über das Konzept des formalen Ableitens auf andere Grundkörper erweitern.


Es sei $K$ ein
\definitionsverweis {Körper}{}{} und sei

\mathl{K[X]}{} der \definitionsverweis {Polynomring}{}{} über $K$. Zu einem Polynom


\mavergleichskette

{\vergleichskette

{ F
}

{ = }{ \sum_{ i = 0 }^{  n   }       a _{ i }  X ^{ i}
}

{ \in }{ K[X]
}

{    }{
}

{   }{
}

}
{}{}{}
heißt das Polynom


\mavergleichskettedisp

{\vergleichskette

{ F'
}

{ =} { n a_nX^{n-1} + (n-1)a_{n-1}X^{n-2} + \cdots + 3a_3X^2 +2a_2X+a_1
}

{ } {
}

{ } {
}

{ } {
}

}
{}{}{}
die \definitionswort {formale Ableitung}{} von $F$.


\inputaufgabe

{}

{


Bestimme die
\definitionsverweis {formale Ableitung}{}{}
von


\mavergleichskettedisp

{\vergleichskette

{ 2X^7+X^6+2X^5+X^4+X^3+X^2+2
}

{ \in} { \Z/( 3 ) [X]
}

{ } {
}

{ } {
}

{ } {
}

}
{}{}{.}


}

{} {}


\inputaufgabe

{}

{


Es sei $K$ ein
\definitionsverweis {Körper}{}{} und sei

\mathl{K[X]}{} der \definitionsverweis {Polynomring}{}{} über $K$. Beweise die folgenden Rechenregeln für das
\definitionsverweis {formale Ableiten}{}{}

\mathl{F \mapsto F'}{:}
\aufzaehlungdrei{Die Ableitung eines konstanten Polynoms ist $0$.
}{Die Ableitung ist
$K$-\definitionsverweis {linear}{}{.}
}{Es gilt die \stichwort {Produktregel} {,} also


\mavergleichskettedisp

{\vergleichskette

{ (FG)'
}

{ =} { FG'+F'G
}

{ } {
}

{ } {
}

{ } {
}

}
{}{}{.}
}


}

{} {}


\inputaufgabe

{}

{


Es sei $K$ ein
\definitionsverweis {Körper}{}{} und sei

\mathl{K[X]}{} der \definitionsverweis {Polynomring}{}{} über $K$. Es sei


\mavergleichskette

{\vergleichskette

{ F
}

{ \in }{ K[X]
}

{  }{
}

{    }{
}

{   }{
}

}
{}{}{}
und


\mavergleichskette

{\vergleichskette

{ a
}

{ \in }{ K
}

{  }{
}

{    }{
}

{   }{
}

}
{}{}{.}
Zeige, dass $a$ genau dann eine
\definitionsverweis {mehrfache Nullstelle}{}{}
von $F$ ist, wenn


\mavergleichskette

{\vergleichskette

{ F'(a)
}

{ = }{ 0
}

{  }{
}

{    }{
}

{   }{
}

}
{}{}{}
ist, wobei $F'$ die
\definitionsverweis {formale Ableitung}{}{}
von $F$ bezeichnet.


}

{} {}


\inputaufgabe

{}

{


Es sei $K$ ein
\definitionsverweis {Körper}{}{}
der
\definitionsverweis {positiven Charakteristik}{}{}


\mavergleichskette

{\vergleichskette

{p
}

{ > }{ 0
}

{  }{
}

{    }{
}

{   }{
}

}
{}{}{.}
Bestimme die Menge der Polynome


\mavergleichskette

{\vergleichskette

{F
}

{ \in }{ K[T]
}

{  }{
}

{    }{
}

{   }{
}

}
{}{}{}
mit
\definitionsverweis {formaler Ableitung}{}{}


\mavergleichskette

{\vergleichskette

{F'
}

{ = }{ 0
}

{  }{
}

{    }{
}

{   }{
}

}
{}{}{.}


}

{} {}


\inputaufgabe

{}

{


Zeige, dass über einem Körper $K$ der
\definitionsverweis {Charakteristik}{}{}
$0$ für das
\definitionsverweis {formale Ableiten}{}{}
die Beziehung
\zusatzklammer {

\mavergleichskettek

{\vergleichskettek

{ i
}

{ \leq }{ n
}

{  }{
}

{    }{
}

{   }{
}

}
{}{}{}} {} {}


\mavergleichskettedisp

{\vergleichskette

{ { \frac{   1  }{ i ! } } \left(  X^n  \right)^{(i)}
}

{ =} { \binom {  n   } { i }  X^{n-i}
}

{ } {
}

{ } {
}

{ } {
}

}
{}{}{}
gilt.


}

{} {}


\inputaufgabe

{}

{


Es sei $K$ ein Körper der
\definitionsverweis {Charakteristik}{}{}
$0$ und sei


\mathbed {f \in K[X]} {}

 {f \neq 0} {}

 {} {} {} {,}
und


\mavergleichskette

{\vergleichskette

{ a
}

{ \in }{ K
}

{  }{
}

{    }{
}

{   }{
}

}
{}{}{.}
Zeige, dass die folgenden \anfuehrung{Ordnungen}{} von $f$ an der Stelle $a$ übereinstimmen.
\aufzaehlungdreiabc{Die Verschwindungsordnung von $f$ an der Stelle $a$, also die minimale Ordnung einer
\definitionsverweis {formalen Ableitung}{}{}
mit


\mavergleichskette

{\vergleichskette

{ f^{(k)}(a)
}

{ \neq }{ 0
}

{  }{
}

{    }{
}

{   }{
}

}
{}{}{.}
}{Der Exponent des Linearfaktors

\mathl{X-a}{} in der Zerlegung von $f$.
}{Die
\definitionsverweis {Ordnung}{}{}
von $f$ an der
\definitionsverweis {Lokalisierung}{}{}


\mathl{K[X]_{(X-a)}}{}  von

\mathl{K[X]}{} am
\definitionsverweis {maximalen Ideal}{}{}


\mathl{(X-a)}{.}
}


}

{} {}


Es sei $R$ ein
\definitionsverweis {kommutativer Ring}{}{}
und


\mavergleichskette

{\vergleichskette

{ {\mathfrak a}
}

{ \subseteq }{ R
}

{  }{
}

{    }{
}

{   }{
}

}
{}{}{}
ein
\definitionsverweis {Ideal}{}{.}
Zu einem
\definitionsverweis {Untermodul}{}{}


\mavergleichskette

{\vergleichskette

{U
}

{ \subseteq }{ V
}

{  }{
}

{    }{
}

{   }{
}

}
{}{}{}
eines
$R$-\definitionsverweis {Moduls}{}{}
$V$ bezeichnet man mit

\mathl{{\mathfrak a} U}{} den von allen Produkten


\mathdisp {fv \text{ mit } f \in {\mathfrak a} \text{ und }v \in U} {  }


\definitionsverweis {erzeugten}{}{}
Untermodul.


\inputaufgabe

{}

{


Es seien


\mavergleichskette

{\vergleichskette

{ {\mathfrak a} , {\mathfrak b}
}

{ \subseteq }{ R
}

{  }{
}

{    }{
}

{   }{
}

}
{}{}{}
\definitionsverweis {Ideale}{}{}
in einem
\definitionsverweis {kommutativen Ring}{}{.}
Zeige, dass das
\definitionsverweis {Idealprodukt}{}{}


\mathl{{\mathfrak a} {\mathfrak b}}{} mit dem
\definitionsverweis {Produkt}{}{}


\mathl{{\mathfrak a} {\mathfrak b}}{} aus dem Ideal ${\mathfrak a}$ und dem
$R$-\definitionsverweis {Untermodul}{}{}


\mavergleichskette

{\vergleichskette

{ {\mathfrak b}
}

{ \subseteq }{ R
}

{  }{
}

{    }{
}

{   }{
}

}
{}{}{}
übereinstimmt.


}

{} {}


\inputaufgabe

{}

{


Es seien


\mavergleichskette

{\vergleichskette

{ {\mathfrak a} , {\mathfrak b}
}

{ \subseteq }{ R
}

{  }{
}

{    }{
}

{   }{
}

}
{}{}{}
\definitionsverweis {Ideale}{}{}
in einem
\definitionsverweis {kommutativen Ring}{}{}
und


\mavergleichskette

{\vergleichskette

{U
}

{ \subseteq }{ V
}

{  }{
}

{    }{
}

{   }{
}

}
{}{}{}
ein
$R$-\definitionsverweis {Untermodul}{}{}
eines
$R$-\definitionsverweis {Moduls}{}{}
$V$. Zeige


\mavergleichskettedisp

{\vergleichskette

{ \left(  {\mathfrak a} \cdot {\mathfrak b}  \right)  \cdot U
}

{ =} { {\mathfrak a} \cdot \left(  {\mathfrak b} \cdot U  \right)
}

{ } {
}

{ } {
}

{ } {
}

}
{}{}{.}


}

{} {}


\inputaufgabe

{}

{


Es seien


\mavergleichskette

{\vergleichskette

{ {\mathfrak a} , {\mathfrak b}
}

{ \subseteq }{ R
}

{  }{
}

{    }{
}

{   }{
}

}
{}{}{}
\definitionsverweis {Ideale}{}{}
in einem
\definitionsverweis {kommutativen Ring}{}{}
und sei


\mavergleichskette

{\vergleichskette

{U
}

{ \subseteq }{ V
}

{  }{
}

{    }{
}

{   }{
}

}
{}{}{}
ein
$R$-\definitionsverweis {Untermodul}{}{}
eines
$R$-\definitionsverweis {Moduls}{}{}
$V$. Zeige


\mavergleichskettedisp

{\vergleichskette

{ \left(  {\mathfrak a}+ {\mathfrak b}  \right)  \cdot U
}

{ =} { {\mathfrak a} \cdot U  + {\mathfrak b} \cdot U
}

{ } {
}

{ } {
}

{ } {
}

}
{}{}{.}


}

{} {}


\inputaufgabe

{}

{


Es sei


\mavergleichskette

{\vergleichskette

{ {\mathfrak a}
}

{ \subseteq }{ R
}

{  }{
}

{    }{
}

{   }{
}

}
{}{}{}
ein
\definitionsverweis {Ideal}{}{}
in einem
\definitionsverweis {kommutativen Ring}{}{}
und seien


\mavergleichskette

{\vergleichskette

{U,W
}

{ \subseteq }{ V
}

{  }{
}

{    }{
}

{   }{
}

}
{}{}{}
$R$-\definitionsverweis {Untermoduln}{}{}
eines
$R$-\definitionsverweis {Moduls}{}{}
$V$. Zeige


\mavergleichskettedisp

{\vergleichskette

{ {\mathfrak a}  \cdot \left(  U +W  \right)
}

{ =} { {\mathfrak a} \cdot U  + {\mathfrak a} \cdot W
}

{ } {
}

{ } {
}

{ } {
}

}
{}{}{.}


}

{} {}


\inputaufgabe

{}

{


Es sei

\mathl{(R, {\mathfrak m} )}{} ein
\definitionsverweis {noetherscher}{}{}
\definitionsverweis {lokaler Ring}{}{.}
Zeige, dass aus


\mavergleichskette

{\vergleichskette

{ {\mathfrak m}^{n+1}
}

{ = }{ {\mathfrak m}^n
}

{  }{
}

{    }{
}

{   }{
}

}
{}{}{}
folgt, dass


\mavergleichskette

{\vergleichskette

{ {\mathfrak m}^n
}

{ = }{ 0
}

{  }{
}

{    }{
}

{   }{
}

}
{}{}{}
ist.


}

{} {}


\inputaufgabe

{}

{


Es sei

\mathl{(R, {\mathfrak m} )}{} ein
\definitionsverweis {lokaler}{}{}
\definitionsverweis {Integritätsbereich}{}{,}
der kein
\definitionsverweis {Körper}{}{}
sei. Es sei $Q$ der
\definitionsverweis {Quotientenkörper}{}{}
von $R$. Zeige


\mavergleichskette

{\vergleichskette

{ {\mathfrak m} Q
}

{ = }{ Q
}

{  }{
}

{    }{
}

{   }{
}

}
{}{}{.}


}

{} {}


  \zwischenueberschrift{Aufgaben zum Abgeben}


\inputaufgabe

{4}

{


Es sei $K$ ein
\definitionsverweis {Körper}{}{}
und $K(T)$ der
\definitionsverweis {Körper der rationalen Funktionen}{}{}
über $K$. Finde einen
\definitionsverweis {diskreten Bewertungsring}{}{}


\mavergleichskette

{\vergleichskette

{ R
}

{ \subseteq }{ K(T)
}

{  }{
}

{    }{
}

{   }{
}

}
{}{}{}
mit


\mavergleichskette

{\vergleichskette

{ Q(R)
}

{ = }{ K(T)
}

{  }{
}

{    }{
}

{   }{
}

}
{}{}{}
und mit


\mavergleichskette

{\vergleichskette

{ R \cap K[T]
}

{ = }{ K
}

{  }{
}

{    }{
}

{   }{
}

}
{}{}{.}


}

{} {}


\inputaufgabe

{4}

{


Es sei $K$ ein
\definitionsverweis {Körper}{}{.}
Eine

\definitionswortenp{Potenzreihe in einer Variablen}{} über $K$ ist ein formaler Ausdruck der Form


\mathdisp {a_0+a_1T+a_2T^2+a_3T^3+ \ldots \text{ mit } a_i \in K} { . }


Es kann hier also unendlich viele von $0$ verschiedene Koeffizienten $a_i$ geben. Definiere eine Ringstruktur auf der Menge aller Potenzreihen, die die Ringstruktur auf dem Polynomring in einer Variablen fortsetzt. Zeige, dass dieser Ring ein
\definitionsverweis {diskreter Bewertungsring}{}{}
ist.


}

{} {}


\inputaufgabe

{4}

{


Es sei $R$ ein
\definitionsverweis {Integritätsbereich}{}{}
mit folgender Eigenschaft: zu je zwei Elementen


\mavergleichskette

{\vergleichskette

{ f,g
}

{ \in }{ R
}

{  }{
}

{    }{
}

{   }{
}

}
{}{}{}
gelte, dass $f$ ein Teiler von $g$ ist oder dass $g$ ein Teiler von $f$ ist. Es sei $R$
\definitionsverweis {noethersch}{}{,}
aber kein
\definitionsverweis {Körper}{}{.}
Zeige, dass $R$ ein
\definitionsverweis {diskreter Bewertungsring}{}{}
ist.


}

{} {}


\inputaufgabe

{3}

{


Zeige, dass in

\mathl{K[X,Y]_{(X,Y)}/ \left(  X^2-Y^3  \right)}{} jedes Ideal durch maximal zwei Erzeuger gegeben ist.


}

{} {}


\inputaufgabe

{3}

{


Man gebe ein Beispiel einer ebenen monomialen Kurve und eines Ideals im zugehörigen lokalen Ring der Singularität, das nicht von zwei Elementen erzeugt werden kann.


}

{} {}
